% GENERATED FILE. Edit canonical lessons, then run npm run generate:books.
% canonical-source-hash: fnv1a64:01deab7ab83a6229
% canonical-lessons: ZH-C235-weile, ZH-C235-youyu, ZH-C235-liyou, ZH-C235-zhihao, ZH-C235-guaibude, ZH-R235-review-230-235

\chapter{For a Reason}
\label{ch:zh-for-a-reason}
\hlchaptermodality{235}

\begin{chapteropening}
\textbf{By the end of this chapter:} \emph{I can give reasons and purposes: what I do it for, what caused it, my grounds, what I had to do, and no wonder.}

Complete the last lesson of chapter 235: I can give reasons and purposes: what I do it for, what caused it, my grounds, what I had to do, and no wonder.
\end{chapteropening}

\section[\texorpdfstring{\zh{为了} (wèile)}{为了 (wèile)}]{\zh{为了} (wèile) — for; in order to}

\label{lesson:ZH-C235-weile}



\begin{quote}
Before the new one: say the Mandarin for a raincoat, then the Mandarin for to shelter from the rain.
\end{quote}

\subsection*{You'll want to know: \zh{为了}}
\textbf{\zh{为了}} — \emph{wèile} — \textquotedblleft{}for; in order to\textquotedblright{}.

\begin{quote}
\textbf{\zh{为了身体我们都不吃点心了}} \emph{wèile shēntǐ wǒmen dōu bù chī diǎnxin le} — For our health, we've all given up snacks.
\end{quote}

\begin{cousinweb}
\textbf{\zh{为了}} (\emph{wèile}) is built from \zh{为} 'for' and \zh{了}, both inherited from Middle Chinese, \emph{hjweH} and \emph{lewX} (Guangyun, in Baxter's transcription), which glosses the first 'to help'; in the word the second loses its tone.
\end{cousinweb}

\begin{grammarlens}[title={What you've built}]
One more word for giving reasons.
\end{grammarlens}

\subsection*{Your turn}
\begin{itemize}
\raggedright
  \item \emph{Say it:} \emph{wèile}
  \item \emph{Say it:} \emph{wèile}, once more
  \item \emph{Say it:} \emph{wèile shēntǐ wǒmen dōu bù chī diǎnxin le} — For our health, we've all given up snacks.
  \item \emph{Recall:} say the Mandarin for a raincoat, then the Mandarin for to shelter from the rain, then say \emph{wèile} again
\end{itemize}

\begin{tcolorbox}[breakable,colback=teal!4,colframe=teal!35!black,title={Before you move on}]
What does \zh{为了} mean? (\textquotedblleft{}For; in order to\textquotedblright{}.) Say it once more.
\end{tcolorbox}
\section[yóuyú]{yóuyú — owing to; because of}

\label{lesson:ZH-C235-youyu}



\begin{quote}
Before the new one: say the Mandarin for to shelter from the rain, then the Mandarin for in order to.
\end{quote}

\subsection*{You'll want to know: yóuyú}
\emph{yóuyú} — \textquotedblleft{}owing to; because of\textquotedblright{}.

\begin{quote}
\emph{yóuyú tiānqì bù hǎo wǒmen méi qù hǎibiān} — Because of the bad weather, we didn't go to the seaside.
\end{quote}

Its characters come later; for now it is a word for the ear.

\begin{cousinweb}
\emph{yóuyú} is built from \emph{yóu} 'from' and \emph{yú} 'at', both inherited from Middle Chinese, \emph{yuw} and \emph{'jo} (Guangyun, in Baxter's transcription), which glosses the first 'to follow' and 'by way of', and the second 'to dwell' and lists it as a particle.
\end{cousinweb}

\begin{grammarlens}[title={What you've built}]
One more word for giving reasons.
\end{grammarlens}

\subsection*{Your turn}
\begin{itemize}
\raggedright
  \item \emph{Say it:} \emph{yóuyú}
  \item \emph{Say it:} \emph{yóuyú}, once more
  \item \emph{Say it:} \emph{yóuyú tiānqì bù hǎo wǒmen méi qù hǎibiān} — Because of the bad weather, we didn't go to the seaside.
  \item \emph{Recall:} say the Mandarin for to shelter from the rain, then the Mandarin for in order to, then say \emph{yóuyú} again
\end{itemize}

\begin{tcolorbox}[breakable,colback=teal!4,colframe=teal!35!black,title={Before you move on}]
What does yóuyú mean? (\textquotedblleft{}Owing to; because of\textquotedblright{}.) Say it once more.
\end{tcolorbox}
\section[lǐyóu]{lǐyóu — reason; grounds}

\label{lesson:ZH-C235-liyou}



\begin{quote}
Before the new one: say the Mandarin for in order to, then the Mandarin for owing to.
\end{quote}

\subsection*{You'll want to know: lǐyóu}
\emph{lǐyóu} — \textquotedblleft{}reason; grounds\textquotedblright{}.

\begin{quote}
\emph{nǐ yǒu shénme lǐyóu} — What are your grounds?
\end{quote}

Its characters come later; for now it is a word for the ear.

\begin{cousinweb}
\emph{lǐyóu} is built from \emph{lǐ} 'reason' and \emph{yóu} 'cause', both inherited from Middle Chinese, \emph{liX} and \emph{yuw} (Guangyun, in Baxter's transcription), which glosses the first 'to put in order' and explains it, after an older book, as working jade.
\end{cousinweb}

\begin{grammarlens}[title={What you've built}]
One more word for giving reasons.
\end{grammarlens}

\subsection*{Your turn}
\begin{itemize}
\raggedright
  \item \emph{Say it:} \emph{lǐyóu}
  \item \emph{Say it:} \emph{lǐyóu}, once more
  \item \emph{Say it:} \emph{nǐ yǒu shénme lǐyóu} — What are your grounds?
  \item \emph{Recall:} say the Mandarin for in order to, then the Mandarin for owing to, then say \emph{lǐyóu} again
\end{itemize}

\begin{tcolorbox}[breakable,colback=teal!4,colframe=teal!35!black,title={Before you move on}]
What does lǐyóu mean? (\textquotedblleft{}Reason; grounds\textquotedblright{}.) Say it once more.
\end{tcolorbox}
\section[zhǐhǎo]{zhǐhǎo — to have no choice but to}

\label{lesson:ZH-C235-zhihao}



\begin{quote}
Before the new one: say the Mandarin for owing to, then the Mandarin for grounds for doing something.
\end{quote}

\subsection*{You'll want to know: zhǐhǎo}
\emph{zhǐhǎo} — \textquotedblleft{}to have no choice but to\textquotedblright{}.

\begin{quote}
\emph{méiyǒu chē le wǒmen zhǐhǎo zǒulù} — There was no transport left, so we had to walk.
\end{quote}

Its characters come later; for now it is a word for the ear.

Tone: zhǐ and hǎo are both third tones, so the first rises and it sounds like zhíhǎo.

\begin{cousinweb}
\emph{zhǐhǎo} is built from \emph{zhǐ} 'only' and \zh{好} 'good'. The second is inherited from Middle Chinese \emph{xawX} (Guangyun, in Baxter's transcription); the dictionary lists the first, \emph{tsyeX} there, as a particle.
\end{cousinweb}

\begin{grammarlens}[title={What you've built}]
One more word for giving reasons.
\end{grammarlens}

\subsection*{Your turn}
\begin{itemize}
\raggedright
  \item \emph{Say it:} \emph{zhǐhǎo}
  \item \emph{Say it:} \emph{zhǐhǎo}, once more
  \item \emph{Say it:} \emph{méiyǒu chē le wǒmen zhǐhǎo zǒulù} — There was no transport left, so we had to walk.
  \item \emph{Recall:} say the Mandarin for owing to, then the Mandarin for grounds for doing something, then say \emph{zhǐhǎo} again
\end{itemize}

\begin{tcolorbox}[breakable,colback=teal!4,colframe=teal!35!black,title={Before you move on}]
What does zhǐhǎo mean? (\textquotedblleft{}To have no choice but to\textquotedblright{}.) Say it once more.
\end{tcolorbox}
\section[guàibude]{guàibude — no wonder}

\label{lesson:ZH-C235-guaibude}



\begin{quote}
Before the new one: say the Mandarin for grounds for doing something, then the Mandarin for to have no choice but.
\end{quote}

\subsection*{You'll want to know: guàibude}
\emph{guàibude} — \textquotedblleft{}no wonder\textquotedblright{}.

\begin{quote}
\emph{tā gǎnmào le guàibude tā méi lái} — He's got a cold; no wonder he didn't come.
\end{quote}

Its characters come later; for now it is a word for the ear.

Tone: \zh{不} is light here, bu, with no tone of its own.

\begin{cousinweb}
\emph{guàibude} is built from \emph{guài} 'to find strange', \zh{不} 'not' and \zh{得} 'can', all inherited from Middle Chinese, \emph{kweajH}, \emph{pjut} and \emph{tok} (Guangyun, in Baxter's transcription), which glosses the first 'strange'; in the word the last two lose their tones.
\end{cousinweb}

\begin{grammarlens}[title={What you've built}]
One more word for giving reasons.
\end{grammarlens}

\subsection*{Your turn}
\begin{itemize}
\raggedright
  \item \emph{Say it:} \emph{guàibude}
  \item \emph{Say it:} \emph{guàibude}, once more
  \item \emph{Say it:} \emph{tā gǎnmào le guàibude tā méi lái} — He's got a cold; no wonder he didn't come.
  \item \emph{Recall:} say the Mandarin for grounds for doing something, then the Mandarin for to have no choice but, then say \emph{guàibude} again
\end{itemize}

\begin{tcolorbox}[breakable,colback=teal!4,colframe=teal!35!black,title={Before you move on}]
What does guàibude mean? (\textquotedblleft{}No wonder\textquotedblright{}.) Say it once more.
\end{tcolorbox}
\section[(dialogue)]{Looking back: from Head to Toe to For a Reason}

\label{lesson:ZH-R235-review-230-235}



\begin{quote}
Say the Mandarin for on your person and body weight.
\end{quote}

\subsection*{Your turn}
\begin{itemize}
\raggedright
  \item \emph{Say it:} on your person, body weight, a tummy, a nose, a leg
  \item \emph{Say it:} Chinese medicine, western medicine, to register at a clinic, a nurse, to prescribe
  \item \emph{Say it:} to have an accident, to get hurt, to bleed, to crash into, an ambulance
  \item \emph{Say it:} appearance, glasses for the eyes, a beard, slender, skin
  \item \emph{Say it:} to put off, to call off, some other day, a raincoat, to shelter from the rain
  \item \emph{Say it:} in order to, owing to, grounds for doing something, to have no choice but, no wonder
  \item \emph{Read it:} \zh{身} \zh{体} \zh{医} \zh{西} — each one aloud, then a word that holds it
\end{itemize}

\begin{tcolorbox}[breakable,colback=teal!4,colframe=teal!35!black,title={Before you move on}]
On your person? (\textbf{\zh{身上}}, \emph{shēnshang}.) Body weight? (\textbf{\zh{体重}}, \emph{tǐzhòng}.) A tummy? (\emph{dùzi}.) A nose? (\emph{bízi}.) A leg? (\emph{tuǐ}.) Chinese medicine? (\textbf{\zh{中医}}, \emph{zhōngyī}.) Western medicine? (\textbf{\zh{西医}}, \emph{xīyī}.) To register at a clinic? (\emph{guàhào}.) A nurse? (\emph{hùshi}.) To prescribe? (\emph{kāiyào}.) To have an accident? (\textbf{\zh{出事}}, \emph{chūshì}.) To get hurt? (\emph{shòushāng}.) To bleed? (\emph{liúxuè}.) To crash into? (\emph{zhuàng}.) An ambulance? (\emph{jiùhùchē}.) Appearance? (\emph{yàngzi}.) Glasses for the eyes? (\emph{yǎnjìng}.) A beard? (\emph{húzi}.) Slender? (\emph{miáotiao}.) Skin? (\emph{pífū}.) To put off? (\emph{tuīchí}.) To call off? (\emph{qǔxiāo}.) Some other day? (\emph{gǎitiān}.) A raincoat? (\textbf{\zh{雨衣}}, \emph{yǔyī}.) To shelter from the rain? (\emph{duǒyǔ}.) In order to? (\textbf{\zh{为了}}, \emph{wèile}.) Owing to? (\emph{yóuyú}.) Grounds for doing something? (\emph{lǐyóu}.) To have no choice but? (\emph{zhǐhǎo}.) No wonder? (\emph{guàibude}.)
\end{tcolorbox}
